The gate is meant to make skill registries harder to game, and the paper describes how they can be gamed: by re-checking results, shopping across configurations, resubmitting, embedding encoded answer tables, or comparing against a stale baseline. These attacks are simple and already possible against registries that publish with-and-without verdicts; describing them alongside the defences should help registries more than attackers. Published transcripts can contain data from task environments, so task classes should use synthetic or licensed data, and redaction should happen before hashing under a registered policy. Logged records are permanent, so no personal data should be written to the log. A certified skill raises an agent's capability on a task class, which can include misuse; the registry's choice of task classes, not the statistics, is where that is controlled.
